% Phase 2 hybrid models and final comparison (Hardik). \input from main.tex.
% Every number comes from results/ files via tox.hybrid_report / tox.gnn_report tables;
% nothing in this section was measured outside tox.metrics.
\section{Hybrid Models}
\label{sec:hybrid}

\subsection{Rationale and the three variants}
The two parents of our Phase 2 comparison make different mistakes. The fingerprint models
encode a fixed vocabulary of substructures, so they are strong on scaffolds they have seen
variations of and weak on genuinely new chemistry; the graph networks learn task-specific
representations from the same training rows, but with 3--16\% positives per endpoint they
overfit the sparse endpoints (Section~\ref{sec:gnn-models}). If the error profiles are
sufficiently different, combining the two should help most on the sparse endpoints and on
the low-similarity test molecules --- and that is a testable claim, not a slogan
(Section~\ref{sec:hybrid-ablation}). We combine the two parents three ways, in increasing
order of coupling:

\begin{description}
\item[hyb\_embxgb] The pooled readout vector of the frozen, tuned MPNN
(Section~\ref{sec:gnn-models}; encoder weights from the same split and seed, so it was
trained on this split's training rows only) is concatenated with the 2048-bit Morgan
fingerprint, and one XGBoost is fitted per endpoint on the concatenation, tuned on
validation PR-AUC exactly like the Phase 1 grid.
\item[hyb\_stack] A per-endpoint logistic regression is fitted on the \emph{validation}
probabilities of \texttt{xgb\_fpdesc} and \texttt{mpnn} and applied to test. The
meta-learner is never fitted on test. Endpoints whose validation labels do not contain
both classes have nothing to learn: for them the stack falls back to the mean of the two
parents.
\item[hyb\_fusion] The joint model: the Morgan fingerprint goes through a small MLP branch,
the graph through the MPNN encoder branch (warm-started from the checkpoint of the same
split and seed), the two vectors are concatenated before one shared head, and the whole
network is trained end to end with the masked weighted BCE of
Section~\ref{sec:gnn-arch}, the Adam recipe of Section~\ref{sec:gnn-models}, and the same
early stopping (validation mean ROC-AUC, patience 25). The graph branch is initialised
from the parent checkpoint, so fusion starts from a trained message passer and fine-tunes
it jointly with the fingerprint branch.
\end{description}

Two rules keep the comparison honest (team plan, p.~19): the GNN encoder is only ever used
with weights trained on the same split's training part --- reusing an encoder trained on
all data would leak; and the stacking meta-learner is fitted on validation, never on test.
Every hybrid is reported against both parents on the same splits and seeds
(Table~\ref{tab:hybrid-main}).

\subsection{Results}
Table~\ref{tab:hybrid-main} is the final comparison on the shared splits and scorer. Three
findings survive the seed spread and the Wilcoxon tests of Section~\ref{sec:validation}:

\begin{enumerate}
\item \textbf{The GNNs close most, but not all, of the gap to the fingerprint models.}
On the randomized scaffold splits the graph models reach 0.813--0.828 mean test ROC-AUC
against the classical models' 0.775--0.839, and all three beat the corresponding Phase 1
fingerprint-only models; but \texttt{rf\_fpdesc} stays ahead of every graph model
(Wilcoxon $p=0.02$--$0.03$, Table~\ref{tab:hybrid-wilcoxon}), in line with what
Yang et al.~\cite{yang2019analyzing} report for well-tuned fingerprint models.
\item \textbf{Stacking is the hybrid that works.} \texttt{hyb\_stack} beats both of its
parents on the scaffold and random split families (its parents score 0.813 and 0.812 mean
ROC-AUC on the scaffold splits; the stack 0.836) and is the only model the Wilcoxon test
cannot separate from \texttt{rf\_fpdesc} ($p=0.47$); on the random splits it is ahead on
paper (0.864 vs 0.862) with the same seed spread.
\item \textbf{End-to-end fusion is a negative result with a clear cause.}
\texttt{hyb\_fusion} reaches only 0.748/0.311 (ROC/PR) on \texttt{scaffold\_det}, below
both parents. Fine-tuning the whole 256-unit encoder at the GNN recipe's learning rate
$10^{-3}$ degrades within a few epochs, so early stopping freezes the epoch-0 or epoch-1
(warm-start) weights in 7 of 9 runs. The ablations locate the failure: removing the fingerprint branch
changes almost nothing (0.744), while removing the graph branch collapses to 0.671 - the
joint head alone cannot learn what the fingerprints know.
\end{enumerate}

The sparse/dense and similarity-bin analyses (Tables~\ref{tab:hybrid-sparse}
and~\ref{tab:hybrid-simbins}) test the stated rationale of the plan: the gap between the
best fingerprint model and the best hybrid is widest on the 4 sparsest endpoints and on
the $[0,0.3)$ Tanimoto bin, where \texttt{hyb\_stack} (0.724) clearly beats its graph
parent \texttt{mpnn} (0.645) - but \texttt{rf\_fpdesc} is also strong there (0.751), so
the honest conclusion is that the stacking gains come from averaging two different error
profiles, not from a special affinity of hybrids for unseen scaffolds.

% GENERATED by python -m tox.hybrid_report from results/ files (no new metrics)

\begin{table*}[t]
\centering
\caption{Final Phase 2 comparison: test ROC-AUC and PR-AUC, 12-endpoint mean, mean $\pm$ std over 3 runs. All models share the splits, the fingerprints and the scorer; the hybrids differ only in how the two parents are combined.}
\label{tab:hybrid-main}
\small
\begin{tabular}{lcccccc}
\toprule
 & \multicolumn{3}{c}{ROC-AUC} & \multicolumn{3}{c}{PR-AUC} \\
model & scaffold\_det & scaffold & random & scaffold\_det & scaffold & random \\
\midrule
lr\_fp & 0.715 $\pm$ 0.000 & 0.775 $\pm$ 0.036 & 0.801 $\pm$ 0.013 & 0.319 $\pm$ 0.000 & 0.395 $\pm$ 0.051 & 0.435 $\pm$ 0.038 \\
lr\_fpdesc & 0.722 $\pm$ 0.000 & 0.788 $\pm$ 0.039 & 0.826 $\pm$ 0.010 & 0.299 $\pm$ 0.000 & 0.367 $\pm$ 0.068 & 0.463 $\pm$ 0.029 \\
rf\_fp & 0.732 $\pm$ 0.001 & 0.803 $\pm$ 0.019 & 0.833 $\pm$ 0.009 & 0.366 $\pm$ 0.001 & 0.442 $\pm$ 0.043 & 0.484 $\pm$ 0.019 \\
rf\_fpdesc & 0.782 $\pm$ 0.003 & 0.839 $\pm$ 0.024 & 0.862 $\pm$ 0.002 & 0.406 $\pm$ 0.002 & 0.481 $\pm$ 0.043 & 0.531 $\pm$ 0.016 \\
xgb\_fp & 0.696 $\pm$ 0.003 & 0.767 $\pm$ 0.045 & 0.801 $\pm$ 0.012 & 0.306 $\pm$ 0.006 & 0.389 $\pm$ 0.047 & 0.435 $\pm$ 0.024 \\
xgb\_fpdesc & 0.761 $\pm$ 0.008 & 0.813 $\pm$ 0.032 & 0.855 $\pm$ 0.007 & 0.391 $\pm$ 0.008 & 0.446 $\pm$ 0.057 & 0.520 $\pm$ 0.020 \\
gcn & 0.752 $\pm$ 0.007 & 0.826 $\pm$ 0.016 & 0.842 $\pm$ 0.008 & 0.354 $\pm$ 0.004 & 0.446 $\pm$ 0.052 & 0.460 $\pm$ 0.023 \\
mpnn & 0.732 $\pm$ 0.011 & 0.813 $\pm$ 0.020 & 0.831 $\pm$ 0.007 & 0.309 $\pm$ 0.014 & 0.386 $\pm$ 0.014 & 0.435 $\pm$ 0.049 \\
afp & 0.769 $\pm$ 0.002 & 0.828 $\pm$ 0.017 & 0.844 $\pm$ 0.010 & 0.354 $\pm$ 0.011 & 0.419 $\pm$ 0.052 & 0.457 $\pm$ 0.009 \\
hyb\_fusion & 0.748 $\pm$ 0.008 & 0.802 $\pm$ 0.024 & 0.833 $\pm$ 0.008 & 0.311 $\pm$ 0.020 & 0.376 $\pm$ 0.035 & 0.475 $\pm$ 0.036 \\
hyb\_embxgb & 0.742 $\pm$ 0.009 & 0.830 $\pm$ 0.019 & 0.843 $\pm$ 0.008 & 0.336 $\pm$ 0.006 & 0.441 $\pm$ 0.048 & 0.495 $\pm$ 0.035 \\
hyb\_stack & 0.768 $\pm$ 0.006 & 0.836 $\pm$ 0.022 & 0.864 $\pm$ 0.002 & 0.396 $\pm$ 0.006 & 0.465 $\pm$ 0.038 & 0.529 $\pm$ 0.022 \\
\midrule
positive rate (random PR-AUC) & \multicolumn{3}{c}{} & 0.103 & 0.103 & 0.077 \\
\bottomrule
\end{tabular}
\end{table*}

\begin{table}[t]
\centering
\caption{Hybrid ablations on \texttt{scaffold\_det}: test ROC-AUC and PR-AUC (12-endpoint mean, mean $\pm$ std over seeds). One factor changes per arm; the baseline is the main run of the same variant.}
\label{tab:hybrid-ablation}
\small
\begin{tabular}{llccc}
\toprule
variant & arm & runs & ROC-AUC & PR-AUC \\
\midrule
hyb\_embxgb & finetune & 3 & 0.749 $\pm$ 0.005 & 0.352 $\pm$ 0.014 \\
hyb\_fusion & nofp & 3 & 0.744 $\pm$ 0.009 & 0.322 $\pm$ 0.010 \\
hyb\_fusion & nograph & 3 & 0.671 $\pm$ 0.005 & 0.266 $\pm$ 0.006 \\
hyb\_embxgb & baseline & 3 & 0.742 $\pm$ 0.009 & 0.336 $\pm$ 0.006 \\
hyb\_fusion & baseline & 3 & 0.748 $\pm$ 0.008 & 0.311 $\pm$ 0.020 \\
hyb\_stack & baseline & 3 & 0.768 $\pm$ 0.006 & 0.396 $\pm$ 0.006 \\
\bottomrule
\end{tabular}
\end{table}

\begin{table}[t]
\centering
\caption{Test ROC-AUC and PR-AUC on the 4 sparsest vs the 4 densest endpoints (\texttt{scaffold\_det}), the plan's direct test of the hybrid rationale: combining two parents should help most where labels are rare.}
\label{tab:hybrid-sparse}
\small
\begin{tabular}{llcc}
\toprule
model & endpoints & ROC-AUC & PR-AUC \\
\midrule
afp & 4 sparsest & 0.764 $\pm$ 0.040 & 0.342 $\pm$ 0.077 \\
afp & 4 densest & 0.778 $\pm$ 0.085 & 0.477 $\pm$ 0.106 \\
gcn & 4 sparsest & 0.740 $\pm$ 0.051 & 0.330 $\pm$ 0.073 \\
gcn & 4 densest & 0.760 $\pm$ 0.081 & 0.474 $\pm$ 0.093 \\
hyb\_embxgb & 4 sparsest & 0.742 $\pm$ 0.043 & 0.295 $\pm$ 0.085 \\
hyb\_embxgb & 4 densest & 0.757 $\pm$ 0.073 & 0.464 $\pm$ 0.087 \\
hyb\_fusion & 4 sparsest & 0.760 $\pm$ 0.034 & 0.284 $\pm$ 0.099 \\
hyb\_fusion & 4 densest & 0.751 $\pm$ 0.063 & 0.430 $\pm$ 0.085 \\
hyb\_stack & 4 sparsest & 0.755 $\pm$ 0.022 & 0.369 $\pm$ 0.059 \\
hyb\_stack & 4 densest & 0.789 $\pm$ 0.081 & 0.536 $\pm$ 0.090 \\
lr\_fp & 4 sparsest & 0.709 $\pm$ 0.022 & 0.335 $\pm$ 0.062 \\
lr\_fp & 4 densest & 0.733 $\pm$ 0.069 & 0.419 $\pm$ 0.097 \\
lr\_fpdesc & 4 sparsest & 0.717 $\pm$ 0.046 & 0.255 $\pm$ 0.075 \\
lr\_fpdesc & 4 densest & 0.735 $\pm$ 0.093 & 0.429 $\pm$ 0.124 \\
mpnn & 4 sparsest & 0.728 $\pm$ 0.040 & 0.270 $\pm$ 0.077 \\
mpnn & 4 densest & 0.745 $\pm$ 0.070 & 0.426 $\pm$ 0.094 \\
rf\_fp & 4 sparsest & 0.765 $\pm$ 0.029 & 0.386 $\pm$ 0.042 \\
rf\_fp & 4 densest & 0.744 $\pm$ 0.052 & 0.469 $\pm$ 0.086 \\
rf\_fpdesc & 4 sparsest & 0.786 $\pm$ 0.046 & 0.385 $\pm$ 0.072 \\
rf\_fpdesc & 4 densest & 0.791 $\pm$ 0.073 & 0.532 $\pm$ 0.084 \\
xgb\_fp & 4 sparsest & 0.708 $\pm$ 0.033 & 0.329 $\pm$ 0.037 \\
xgb\_fp & 4 densest & 0.712 $\pm$ 0.057 & 0.404 $\pm$ 0.061 \\
xgb\_fpdesc & 4 sparsest & 0.754 $\pm$ 0.078 & 0.373 $\pm$ 0.087 \\
xgb\_fpdesc & 4 densest & 0.783 $\pm$ 0.082 & 0.528 $\pm$ 0.093 \\
\bottomrule
\end{tabular}
\end{table}

\begin{table}[t]
\centering
\caption{Test ROC-AUC (all 12 tasks pooled, missing labels masked) by maximum Tanimoto similarity of each test molecule to the training set (\texttt{scaffold\_det}, mean $\pm$ std over seeds). Low-similarity bins are where unseen scaffolds live.}
\label{tab:hybrid-simbins}
\small
\begin{tabular}{llcc}
\toprule
model & bin & mols & ROC-AUC \\
\midrule
afp & [0, 0.3) & 172 & 0.703 $\pm$ 0.007 \\
afp & [0.3, 0.5) & 365 & 0.798 $\pm$ 0.008 \\
afp & [0.5, 1) & 235 & 0.846 $\pm$ 0.007 \\
gcn & [0, 0.3) & 172 & 0.703 $\pm$ 0.009 \\
gcn & [0.3, 0.5) & 365 & 0.787 $\pm$ 0.007 \\
gcn & [0.5, 1) & 235 & 0.846 $\pm$ 0.005 \\
hyb\_embxgb & [0, 0.3) & 172 & 0.660 $\pm$ 0.036 \\
hyb\_embxgb & [0.3, 0.5) & 365 & 0.770 $\pm$ 0.006 \\
hyb\_embxgb & [0.5, 1) & 235 & 0.835 $\pm$ 0.009 \\
hyb\_fusion & [0, 0.3) & 172 & 0.657 $\pm$ 0.032 \\
hyb\_fusion & [0.3, 0.5) & 365 & 0.766 $\pm$ 0.021 \\
hyb\_fusion & [0.5, 1) & 235 & 0.812 $\pm$ 0.033 \\
hyb\_stack & [0, 0.3) & 172 & 0.724 $\pm$ 0.019 \\
hyb\_stack & [0.3, 0.5) & 365 & 0.787 $\pm$ 0.007 \\
hyb\_stack & [0.5, 1) & 235 & 0.840 $\pm$ 0.006 \\
lr\_fp & [0, 0.3) & 172 & 0.662 $\pm$ 0.000 \\
lr\_fp & [0.3, 0.5) & 365 & 0.723 $\pm$ 0.000 \\
lr\_fp & [0.5, 1) & 235 & 0.819 $\pm$ 0.000 \\
lr\_fpdesc & [0, 0.3) & 172 & 0.704 $\pm$ 0.000 \\
lr\_fpdesc & [0.3, 0.5) & 365 & 0.740 $\pm$ 0.000 \\
lr\_fpdesc & [0.5, 1) & 235 & 0.816 $\pm$ 0.000 \\
mpnn & [0, 0.3) & 172 & 0.645 $\pm$ 0.007 \\
mpnn & [0.3, 0.5) & 365 & 0.769 $\pm$ 0.015 \\
mpnn & [0.5, 1) & 235 & 0.828 $\pm$ 0.014 \\
rf\_fp & [0, 0.3) & 172 & 0.654 $\pm$ 0.013 \\
rf\_fp & [0.3, 0.5) & 365 & 0.723 $\pm$ 0.006 \\
rf\_fp & [0.5, 1) & 235 & 0.823 $\pm$ 0.001 \\
rf\_fpdesc & [0, 0.3) & 172 & 0.751 $\pm$ 0.009 \\
rf\_fpdesc & [0.3, 0.5) & 365 & 0.810 $\pm$ 0.005 \\
rf\_fpdesc & [0.5, 1) & 235 & 0.865 $\pm$ 0.002 \\
xgb\_fp & [0, 0.3) & 172 & 0.639 $\pm$ 0.017 \\
xgb\_fp & [0.3, 0.5) & 365 & 0.694 $\pm$ 0.026 \\
xgb\_fp & [0.5, 1) & 235 & 0.797 $\pm$ 0.013 \\
xgb\_fpdesc & [0, 0.3) & 172 & 0.749 $\pm$ 0.014 \\
xgb\_fpdesc & [0.3, 0.5) & 365 & 0.773 $\pm$ 0.022 \\
xgb\_fpdesc & [0.5, 1) & 235 & 0.829 $\pm$ 0.013 \\
\bottomrule
\end{tabular}
\end{table}

\begin{table}[t]
\centering
\caption{Wilcoxon signed-rank tests on the randomized scaffold splits (12 paired endpoint means; each endpoint averaged over its 3 runs). $^{*}$/$^{**}$: $p<0.05$ / $p<0.01$. Exploratory: 12 endpoints and many pairs, no multiplicity correction.}
\label{tab:hybrid-wilcoxon}
\small
\begin{tabular}{llccc}
\toprule
model $a$ & model $b$ & ROC-AUC $a$ & ROC-AUC $b$ & $p$ \\
\midrule
afp & hyb\_embxgb & 0.769 & 0.742 & 0.007$^{**}$ \\
afp & hyb\_fusion & 0.769 & 0.748 & 0.027$^{*}$ \\
afp & hyb\_stack & 0.769 & 0.768 & 0.970 \\
afp & mpnn & 0.769 & 0.732 & 0.000$^{**}$ \\
afp & rf\_fpdesc & 0.769 & 0.782 & 0.424 \\
afp & xgb\_fpdesc & 0.769 & 0.761 & 0.470 \\
gcn & hyb\_embxgb & 0.752 & 0.742 & 0.301 \\
gcn & hyb\_fusion & 0.752 & 0.748 & 0.910 \\
gcn & hyb\_stack & 0.752 & 0.768 & 0.064 \\
gcn & mpnn & 0.752 & 0.732 & 0.016$^{*}$ \\
gcn & rf\_fpdesc & 0.752 & 0.782 & 0.027$^{*}$ \\
gcn & xgb\_fpdesc & 0.752 & 0.761 & 0.424 \\
hyb\_embxgb & hyb\_fusion & 0.742 & 0.748 & 0.519 \\
hyb\_embxgb & hyb\_stack & 0.742 & 0.768 & 0.002$^{**}$ \\
hyb\_embxgb & mpnn & 0.742 & 0.732 & 0.042$^{*}$ \\
hyb\_embxgb & rf\_fpdesc & 0.742 & 0.782 & 0.002$^{**}$ \\
hyb\_embxgb & xgb\_fpdesc & 0.742 & 0.761 & 0.151 \\
hyb\_fusion & hyb\_stack & 0.748 & 0.768 & 0.042$^{*}$ \\
hyb\_fusion & mpnn & 0.748 & 0.732 & 0.012$^{*}$ \\
hyb\_fusion & rf\_fpdesc & 0.748 & 0.782 & 0.003$^{**}$ \\
hyb\_fusion & xgb\_fpdesc & 0.748 & 0.761 & 0.470 \\
hyb\_stack & mpnn & 0.768 & 0.732 & 0.000$^{**}$ \\
hyb\_stack & rf\_fpdesc & 0.768 & 0.782 & 0.064 \\
hyb\_stack & xgb\_fpdesc & 0.768 & 0.761 & 0.064 \\
lr\_fp & mpnn & 0.715 & 0.732 & 0.151 \\
lr\_fp & rf\_fpdesc & 0.715 & 0.782 & 0.000$^{**}$ \\
lr\_fp & xgb\_fpdesc & 0.715 & 0.761 & 0.002$^{**}$ \\
lr\_fpdesc & mpnn & 0.722 & 0.732 & 0.424 \\
lr\_fpdesc & rf\_fpdesc & 0.722 & 0.782 & 0.001$^{**}$ \\
lr\_fpdesc & xgb\_fpdesc & 0.722 & 0.761 & 0.016$^{*}$ \\
mpnn & rf\_fpdesc & 0.732 & 0.782 & 0.001$^{**}$ \\
mpnn & xgb\_fpdesc & 0.732 & 0.761 & 0.064 \\
rf\_fp & rf\_fpdesc & 0.732 & 0.782 & 0.005$^{**}$ \\
rf\_fp & xgb\_fpdesc & 0.732 & 0.761 & 0.151 \\
rf\_fpdesc & xgb\_fpdesc & 0.782 & 0.761 & 0.012$^{*}$ \\
xgb\_fp & xgb\_fpdesc & 0.696 & 0.761 & 0.001$^{**}$ \\
afp & hyb\_embxgb & 0.828 & 0.830 & 0.733 \\
afp & hyb\_fusion & 0.828 & 0.802 & 0.005$^{**}$ \\
afp & hyb\_stack & 0.828 & 0.836 & 0.077 \\
afp & mpnn & 0.828 & 0.813 & 0.016$^{*}$ \\
afp & rf\_fpdesc & 0.828 & 0.839 & 0.034$^{*}$ \\
afp & xgb\_fpdesc & 0.828 & 0.813 & 0.110 \\
gcn & hyb\_embxgb & 0.826 & 0.830 & 0.380 \\
gcn & hyb\_fusion & 0.826 & 0.802 & 0.007$^{**}$ \\
gcn & hyb\_stack & 0.826 & 0.836 & 0.052 \\
gcn & mpnn & 0.826 & 0.813 & 0.064 \\
gcn & rf\_fpdesc & 0.826 & 0.839 & 0.021$^{*}$ \\
gcn & xgb\_fpdesc & 0.826 & 0.813 & 0.233 \\
hyb\_embxgb & hyb\_fusion & 0.830 & 0.802 & 0.000$^{**}$ \\
hyb\_embxgb & hyb\_stack & 0.830 & 0.836 & 0.129 \\
hyb\_embxgb & lr\_fp & 0.830 & 0.775 & 0.000$^{**}$ \\
hyb\_embxgb & lr\_fpdesc & 0.830 & 0.788 & 0.002$^{**}$ \\
hyb\_embxgb & mpnn & 0.830 & 0.813 & 0.005$^{**}$ \\
hyb\_embxgb & rf\_fp & 0.830 & 0.803 & 0.005$^{**}$ \\
hyb\_embxgb & rf\_fpdesc & 0.830 & 0.839 & 0.151 \\
hyb\_embxgb & xgb\_fp & 0.830 & 0.767 & 0.000$^{**}$ \\
hyb\_embxgb & xgb\_fpdesc & 0.830 & 0.813 & 0.077 \\
hyb\_fusion & hyb\_stack & 0.802 & 0.836 & 0.001$^{**}$ \\
hyb\_fusion & lr\_fp & 0.802 & 0.775 & 0.005$^{**}$ \\
hyb\_fusion & lr\_fpdesc & 0.802 & 0.788 & 0.129 \\
hyb\_fusion & mpnn & 0.802 & 0.813 & 0.077 \\
hyb\_fusion & rf\_fp & 0.802 & 0.803 & 0.910 \\
hyb\_fusion & rf\_fpdesc & 0.802 & 0.839 & 0.001$^{**}$ \\
hyb\_fusion & xgb\_fp & 0.802 & 0.767 & 0.002$^{**}$ \\
hyb\_fusion & xgb\_fpdesc & 0.802 & 0.813 & 0.266 \\
hyb\_stack & lr\_fp & 0.836 & 0.775 & 0.000$^{**}$ \\
hyb\_stack & lr\_fpdesc & 0.836 & 0.788 & 0.001$^{**}$ \\
hyb\_stack & mpnn & 0.836 & 0.813 & 0.000$^{**}$ \\
hyb\_stack & rf\_fp & 0.836 & 0.803 & 0.000$^{**}$ \\
hyb\_stack & rf\_fpdesc & 0.836 & 0.839 & 0.470 \\
hyb\_stack & xgb\_fp & 0.836 & 0.767 & 0.000$^{**}$ \\
hyb\_stack & xgb\_fpdesc & 0.836 & 0.813 & 0.005$^{**}$ \\
lr\_fp & mpnn & 0.775 & 0.813 & 0.001$^{**}$ \\
lr\_fp & rf\_fpdesc & 0.775 & 0.839 & 0.000$^{**}$ \\
lr\_fp & xgb\_fpdesc & 0.775 & 0.813 & 0.001$^{**}$ \\
lr\_fpdesc & mpnn & 0.788 & 0.813 & 0.052 \\
lr\_fpdesc & rf\_fpdesc & 0.788 & 0.839 & 0.001$^{**}$ \\
lr\_fpdesc & xgb\_fpdesc & 0.788 & 0.813 & 0.009$^{**}$ \\
mpnn & rf\_fp & 0.813 & 0.803 & 0.233 \\
mpnn & rf\_fpdesc & 0.813 & 0.839 & 0.002$^{**}$ \\
mpnn & xgb\_fp & 0.813 & 0.767 & 0.000$^{**}$ \\
mpnn & xgb\_fpdesc & 0.813 & 0.813 & 0.850 \\
rf\_fp & rf\_fpdesc & 0.803 & 0.839 & 0.001$^{**}$ \\
rf\_fp & xgb\_fpdesc & 0.803 & 0.813 & 0.339 \\
rf\_fpdesc & xgb\_fp & 0.839 & 0.767 & 0.000$^{**}$ \\
rf\_fpdesc & xgb\_fpdesc & 0.839 & 0.813 & 0.002$^{**}$ \\
xgb\_fp & xgb\_fpdesc & 0.767 & 0.813 & 0.002$^{**}$ \\
afp & hyb\_embxgb & 0.844 & 0.843 & 0.970 \\
afp & hyb\_fusion & 0.844 & 0.833 & 0.027$^{*}$ \\
afp & hyb\_stack & 0.844 & 0.864 & 0.007$^{**}$ \\
afp & mpnn & 0.844 & 0.831 & 0.027$^{*}$ \\
afp & rf\_fpdesc & 0.844 & 0.862 & 0.001$^{**}$ \\
afp & xgb\_fpdesc & 0.844 & 0.855 & 0.077 \\
gcn & hyb\_embxgb & 0.842 & 0.843 & 0.622 \\
gcn & hyb\_fusion & 0.842 & 0.833 & 0.301 \\
gcn & hyb\_stack & 0.842 & 0.864 & 0.003$^{**}$ \\
gcn & mpnn & 0.842 & 0.831 & 0.077 \\
gcn & rf\_fpdesc & 0.842 & 0.862 & 0.016$^{*}$ \\
gcn & xgb\_fpdesc & 0.842 & 0.855 & 0.052 \\
hyb\_embxgb & hyb\_fusion & 0.843 & 0.833 & 0.110 \\
hyb\_embxgb & hyb\_stack & 0.843 & 0.864 & 0.000$^{**}$ \\
hyb\_embxgb & mpnn & 0.843 & 0.831 & 0.064 \\
hyb\_embxgb & rf\_fpdesc & 0.843 & 0.862 & 0.027$^{*}$ \\
hyb\_embxgb & xgb\_fpdesc & 0.843 & 0.855 & 0.052 \\
hyb\_fusion & hyb\_stack & 0.833 & 0.864 & 0.000$^{**}$ \\
hyb\_fusion & mpnn & 0.833 & 0.831 & 0.380 \\
hyb\_fusion & rf\_fpdesc & 0.833 & 0.862 & 0.000$^{**}$ \\
hyb\_fusion & xgb\_fpdesc & 0.833 & 0.855 & 0.003$^{**}$ \\
hyb\_stack & mpnn & 0.864 & 0.831 & 0.000$^{**}$ \\
hyb\_stack & rf\_fpdesc & 0.864 & 0.862 & 0.910 \\
hyb\_stack & xgb\_fpdesc & 0.864 & 0.855 & 0.233 \\
lr\_fp & mpnn & 0.801 & 0.831 & 0.000$^{**}$ \\
lr\_fp & rf\_fpdesc & 0.801 & 0.862 & 0.000$^{**}$ \\
lr\_fp & xgb\_fpdesc & 0.801 & 0.855 & 0.000$^{**}$ \\
lr\_fpdesc & mpnn & 0.826 & 0.831 & 0.970 \\
lr\_fpdesc & rf\_fpdesc & 0.826 & 0.862 & 0.001$^{**}$ \\
lr\_fpdesc & xgb\_fpdesc & 0.826 & 0.855 & 0.001$^{**}$ \\
mpnn & rf\_fpdesc & 0.831 & 0.862 & 0.002$^{**}$ \\
mpnn & xgb\_fpdesc & 0.831 & 0.855 & 0.007$^{**}$ \\
rf\_fp & rf\_fpdesc & 0.833 & 0.862 & 0.001$^{**}$ \\
rf\_fp & xgb\_fpdesc & 0.833 & 0.855 & 0.007$^{**}$ \\
rf\_fpdesc & xgb\_fpdesc & 0.862 & 0.855 & 0.129 \\
xgb\_fp & xgb\_fpdesc & 0.801 & 0.855 & 0.001$^{**}$ \\
\bottomrule
\end{tabular}
\end{table}

\begin{table}[t]
\centering
\caption{Test calibration over all labelled (task, molecule) pairs of every run: expected calibration error (10 equal-width bins) and Brier score.}
\label{tab:hybrid-calibration}
\small
\begin{tabular}{lccc}
\toprule
model & pairs & ECE & Brier \\
\midrule
afp & 65,815 & 0.236 & 0.162 \\
gcn & 65,815 & 0.166 & 0.127 \\
hyb\_embxgb & 65,815 & 0.128 & 0.102 \\
hyb\_fusion & 65,815 & 0.235 & 0.157 \\
hyb\_stack & 65,815 & 0.278 & 0.153 \\
lr\_fp & 65,815 & 0.227 & 0.136 \\
lr\_fpdesc & 65,815 & 0.136 & 0.113 \\
mpnn & 65,815 & 0.226 & 0.155 \\
rf\_fp & 65,815 & 0.066 & 0.076 \\
rf\_fpdesc & 65,815 & 0.025 & 0.061 \\
xgb\_fp & 65,815 & 0.254 & 0.147 \\
xgb\_fpdesc & 65,815 & 0.160 & 0.105 \\
\bottomrule
\end{tabular}
\end{table}


Aggregate scores can hide endpoint-specific differences. Table~\ref{tab:hybrid-endpoint}
therefore reports per-endpoint ROC-AUC and PR-AUC for the fingerprint-plus-descriptor
baseline, the strongest GNN, and the stacking hybrid on the balanced scaffold split family.
No one model is best on every endpoint, and PR-AUC is interpreted relative to each
endpoint's prevalence.

\subsection{Ablations}
\label{sec:hybrid-ablation}
Five ablations ask what each part contributes (plan p.~19); all run on
\texttt{scaffold\_det} with seeds 0--2 and change one factor at a time relative to the
main run of the same variant (Table~\ref{tab:hybrid-ablation}):

\begin{enumerate}
\item \textbf{fusion without the fingerprint branch} (\texttt{hyb\_fusion-nofp}): does the
graph alone, trained with the fusion head, match the fingerprint-informed model?
\item \textbf{fusion without the graph branch} (\texttt{hyb\_fusion-nograph}): does the
fingerprint alone, under the identical protocol, match the joint model?
\item \textbf{frozen vs fine-tuned encoder in hyb\_embxgb} (\texttt{hyb\_embxgb-finetune}):
the encoder is fine-tuned for 30 further epochs at learning rate $10^{-4}$ before its
embeddings are frozen into the XGBoost features. Is the transferred representation already
good, or does it need adaptation?
\item \textbf{sparsest vs densest endpoints}: the hybrid is evaluated on the 4 endpoints
with the fewest labelled positives and on the 4 densest
(Table~\ref{tab:hybrid-sparse}) --- the direct test of the stated rationale.
\item \textbf{hybrid by Tanimoto similarity bin}: each test molecule's maximum Tanimoto
similarity (Morgan 2048) to the training set is binned into $[0,0.3)$, $[0.3,0.5)$,
$[0.5,1.0]$, and ROC-AUC is computed per bin over pooled tasks
(Table~\ref{tab:hybrid-simbins}) --- the direct test of whether the gain comes from the
hard, dissimilar molecules.
\end{enumerate}

If the hybrid loses to its parents, that is reported as a negative result together with the
ablation arms that explain it, in line with what Yang et al.~\cite{yang2019analyzing}
found for well-tuned fingerprint models.
